\chapter{Introduction}
\label{ch:intro}

Every online service that people rely on, whether a shopping site, a hospital
booking portal, or a university registration system, is underneath a pool of
identical server processes sitting behind a single front door. A piece of
software called a \emph{load balancer} stands at that door and decides which
server handles each incoming request. When the pool is the right size and
traffic is steady, this work is invisible and dull, exactly as it should be.

The trouble starts when demand is hard to predict. A sale launches, a news story
breaks, or an exam-results page opens at nine o'clock sharp. Traffic surges, the
fixed pool cannot keep up, requests queue, and the people behind the screens see
the result as slowness or an error page. The usual answer, adding more servers,
often arrives \emph{after} the surge has already done its damage, because the
machinery that detects overload reacts only once latency has climbed.

This report presents \emph{SmartLoad}, a self-hosted middleware that sits in that
doorway and acts \emph{ahead} of the surge instead of after it. SmartLoad watches
the live behaviour of the backend pool, forecasts where demand is heading,
decides how to route requests and whether to grow or shrink the pool, and carries
those decisions out on a standard NGINX \cite{nginx} data plane.

Equally important, every learned decision sits on top of a simple, predictable
rule that takes over the moment the learned part is unavailable or an operator
turns it off. The design follows a single principle: \emph{routing learns; safety
does not}. This chapter explains why the problem matters, what makes it hard,
what we did and did not set out to do, how the method works in plain terms, what
the prototype achieved, and how the rest of the report is organised.

\section{Motivation}
\label{sec:intro:motivation}

The case for SmartLoad spans four impact contexts: economic, societal,
environmental, and cultural. We take each in turn, because a system that sits on
the request path of a production service is never a purely technical thing. It
has consequences for budgets, for the people who depend on the service, for
energy use, and for the operators whose trust it has to earn.

\subsection{Economic and Global Context}
\label{sec:intro:economic}

Load balancing and autoscaling are not exotic functions. They are a routine line
item in the cost of running any networked application. The market for products
that perform them (application-delivery controllers, reverse proxies, and cloud
autoscalers) is mature and crowded.

Open-source data planes such as NGINX \cite{nginx}, HAProxy \cite{haproxy}, and
Envoy \cite{envoy} move traffic for a large share of the public web, while
managed platforms add elastic scaling on top. The Kubernetes Horizontal Pod
Autoscaler \cite{k8shpa} and event-driven autoscalers such as KEDA \cite{keda}
are now standard tooling for teams that have adopted container orchestration.
Cloud infrastructure spending, of which application delivery and elastic compute
are a large part, has grown into one of the biggest categories of enterprise
technology cost, and research on automatic resource provisioning has tracked that
growth for over a decade \cite{lorido2014autoscaling,roy2011autoscaling}.

The tooling that makes elastic scaling truly \emph{proactive}, however, usually
assumes a full orchestration platform and a specialised team to run it.
SmartLoad's intended customer is the team \emph{below} that line: a small or
medium engineering group running a pool of backend services on a handful of
hosts, who needs adaptive traffic management but cannot justify standing up and
staffing a Kubernetes control plane.

For that customer the choices today are poor. One option is to over-provision
permanently, paying for idle capacity to absorb peaks that rarely come. The other
is to wire together a load balancer, a separate autoscaler, a monitoring stack,
and an alerting tool by hand, with no shared decision-making between them.
SmartLoad packages those concerns into a single self-hosted unit that installs
from a container manifest, and that is the economic gap it is built to fill. A
managed control plane is a plausible later commercial direction, but this work is
deliberately scoped as middleware a team can run itself.

\subsection{Societal and Environmental Context}
\label{sec:intro:society}

The reliability of online services is no longer just a convenience question.
People book medical appointments, sit examinations, file taxes, and reach
emergency information through web applications, and an outage or a long slowdown
during a demand peak is felt as a real denial of service to real people. A system
that grows capacity \emph{before} a predictable peak, rather than scrambling
after the alarm, directly improves the experience of everyone waiting behind that
front door, and it helps most during the high-load moments that matter.

The same mechanism brings an environmental benefit. The common way to survive
demand spikes is to leave spare capacity permanently switched on, which means
racks of servers drawing power and cooling to serve a load that arrives only now
and then. Right-sizing compute, scaling the pool \emph{down} when demand recedes
and not only up when it rises, returns that idle capacity, and capacity that is
switched off is energy that is not wasted.

SmartLoad's autoscaler is bidirectional by design for this reason. In our
prototype run it grew the backend pool under a forecast burst and then shrank it
again as the burst subsided, rather than leaving the enlarged pool running. At
the scale of one small deployment the saving is modest. As a default behaviour
across many deployments, scaling idle capacity down becomes a meaningful lever on
the energy footprint of everyday computing.

\subsection{Cultural Context}
\label{sec:intro:culture}

Automation that sits in the request path asks operators to hand over control of
something they are accountable for, and the appetite for that varies widely:
between individuals, between teams, and across the engineering cultures of
different organisations. Some teams will give routing decisions to a learned
model the moment it shows value. Others, for reasons of regulation, risk
tolerance, or hard-won caution, will not trust an opaque component anywhere near
production traffic until they can inspect and reverse every decision it makes. A
system that demands a single fixed level of trust will be rejected by one of these
audiences whichever level it picks.

SmartLoad's answer is to make the level of autonomy a setting rather than an
assumption. Every learned decision is logged and attributable, every learned
component has a deterministic rule behind it, and a single operator flag
(\texttt{safe\_mode}) can switch off all learned behaviour and return the system
to classical routing without taking it offline.

The learned router, in particular, ships pinned to a passive \emph{shadow} mode:
it observes and records what it \emph{would} do, but an operator must explicitly
promote it before it touches live traffic. A kill switch on every decision does
more than guard against failure. It also matches the differing cultural attitudes
toward delegating authority to software, so cautious and confident operators can
adopt the same product, each at the level of autonomy they are prepared to
defend.

\section{Overview of the Key Challenges}
\label{sec:intro:challenges}

Adaptive traffic management is hard for reasons that build on one another, and
naming them up front explains the design choices in the chapters that follow.

The first challenge is timing. Conventional autoscaling is \emph{reactive}: it
notices that latency or utilisation has crossed a threshold and only then adds
capacity. By the time new servers are provisioned and warm, the spike that
triggered the alarm has often already hurt service. Acting ahead of the spike
calls for \emph{forecasting} demand, which the autoscaling literature identifies
as the lever that separates proactive from reactive provisioning
\cite{lorido2014autoscaling,roy2011autoscaling}. Forecasting brings its own
difficulty, because a wrong prediction can scale the pool in the wrong direction.

The second challenge is heterogeneity. Real backend pools are not uniform.
Instances differ in hardware, in cache warmth, and in their position on a shared
host, and a single slow or sick instance can drag down the latency of any request
unlucky enough to land on it. A router that treats all backends as
interchangeable, as plain round-robin does, cannot use these differences. A
router that tries to learn them must be trained on the variety it will face in
production, and training data that lacks that variety produces a model that has
learned little of use.

The third challenge is the risk of putting an opaque model in the request path at
all. A learned routing policy that misbehaves does not produce a wrong answer in
a report; it sends live user traffic to the wrong place. The difficulty is
fundamental: the very adaptiveness that makes a learned component valuable also
makes it risky, and the only way to earn the right to run it in production is to
guarantee that any failure degrades to a safe, predictable baseline rather than
to an outage. This insistence on graceful degradation is a recurring theme in the
autonomic-computing literature \cite{kephart2003autonomic}.

The remaining challenges are infrastructural but no less real. Getting
trustworthy telemetry out of the data plane and into the decision components (the
plumbing of access logs, metrics, and time-series storage) is a real engineering
effort in its own right, and any error in that pipeline quietly corrupts every
downstream decision. Reconfiguring a live load balancer safely, without dropping
in-flight requests or briefly emptying the pool, is a delicate operation that we
designed for carefully so the service stays up during every change.

\section{Project Scope}
\label{sec:intro:scope}

SmartLoad is scoped as a \emph{single-tenant, self-hosted adaptive middleware} for
the request path of a backend pool. Within that scope it handles, as one
integrated whole, the concerns that are usually bought as separate tools: anomaly
detection over backend health, short-horizon demand forecasting, learned request
routing, proactive and bidirectional autoscaling, a central policy and control
surface for operators, and the telemetry pipeline that feeds all of them.

The defining commitment across all of these is that each learned component sits
behind a deterministic fallback and an operator override. The loss of
intelligence reduces the system to classical routing, never to an outage.

Equally important is what SmartLoad deliberately leaves out of this phase. The
prototype is single-tenant: it assumes one operating team on a trusted internal
network. Multi-tenant isolation, role-based access control, operator
authentication, TLS termination at the load balancer, per-token rate limiting,
and an outbound webhook dispatcher are all planned for a future Phase~2 aimed at a
managed, multi-customer offering. These are scope decisions, and they sit outside
the present phase so the work can demonstrate the core adaptive control loop end
to end rather than spreading thinly across a wide feature set.

Several assumptions bound the work. We assume HTTP/1.1 reverse proxying without
session affinity, so the request path carries no hidden per-client state and any
backend can serve any request. We assume the control bus may lose messages, since
it is treated as best-effort, in-memory transport, so every consumer of a learned
signal can tolerate its absence. We assume backends are interchangeable members
of a pool rather than functionally distinct services. And we assume the operator,
not the system, is the final authority: SmartLoad recommends and acts within a
policy, but a human can always inspect, override, or disable any decision.

\section{Proposed Methodology}
\label{sec:intro:method}

In plain terms, SmartLoad runs a continuous \emph{control loop} with four steps
that repeat several times a minute. It \textbf{observes} the live behaviour of
the backend pool: how fast each instance is responding, how many requests are in
flight, and whether any instance is throwing errors. From those observations it
\textbf{predicts} where demand is heading over the next short horizon. Using the
prediction and the current health picture, it \textbf{decides} two things: how to
spread incoming requests across the healthy instances, and whether the pool
should grow, shrink, or hold. Finally it \textbf{acts}, applying those decisions
to the NGINX data plane that carries the traffic, and the loop begins again by
observing the result.

This loop of observe, predict, decide, and act is the classical shape of a
self-managing system, often written MAPE-K (Monitor, Analyse, Plan, Execute over
shared Knowledge), introduced in the autonomic-computing literature
\cite{kephart2003autonomic} that frames our approach.

What makes SmartLoad's version of this loop distinctive is that the decide step is
not a single algorithm but a small committee, applied in a fixed safety order. An
anomaly check runs first and removes any unhealthy backend from consideration,
and that exclusion is never overridden. A learned routing policy then proposes how
to spread traffic across what remains, while a demand forecast drives the scaling
decision in parallel.

Behind each learned component stands a deterministic baseline: a simple threshold
rule for anomalies, a moving average for the forecast, and plain round-robin for
routing. If a learned component is turned off, cannot load, or is overridden by
the operator, its baseline quietly takes its place. The operator sits above the
whole loop: every decision is recorded and attributable, and a single safe-mode
flag collapses the committee back to pure classical routing.

The methodology, then, is not to replace classical infrastructure with a model,
but to wrap classical infrastructure in a learning loop that can always fall back
to it. That fallback is what lets the learned parts be trusted in the request path
at all. \Cref{fig:intro:loop} sketches the loop at a glance; the detailed
architecture is the subject of \cref{ch:design_solution}.

\begin{figure}[t]
\centering
\begin{tikzpicture}[node distance=10mm and 16mm]
  \node[slnode] (obs) {\textbf{Observe}\\telemetry from\\the pool};
  \node[slnode,right=of obs] (pred) {\textbf{Predict}\\demand over\\the next horizon};
  \node[slnode,right=of pred] (dec) {\textbf{Decide}\\route + scale\\(learned, with\\deterministic\\fallback)};
  \node[slnode,right=of dec] (act) {\textbf{Act}\\apply to the\\NGINX data plane};
  \draw[slflow] (obs) -- (pred);
  \draw[slflow] (pred) -- (dec);
  \draw[slflow] (dec) -- (act);
  % closing the loop: act feeds back into the next observation
  \draw[slbus] (act.south) -- ++(0,-9mm) -| node[pos=0.25,below,font=\footnotesize]{feedback: observe the result} (obs.south);
\end{tikzpicture}
\caption{The SmartLoad control loop. Telemetry is observed, demand is predicted, routing and scaling are decided by learned components each backed by a deterministic fallback, and the decisions are applied to the NGINX data plane. The act step feeds the next observation, closing a MAPE-K-style loop.}
\label{fig:intro:loop}
\end{figure}

\section{Summary of the Implementation and Results}
\label{sec:intro:results}

The prototype is a containerised stack of roughly twelve cooperating services
orchestrated with Docker Compose \cite{merkel2014docker}. An NGINX data plane
\cite{nginx} carries the traffic; a telemetry pipeline built on OpenTelemetry
\cite{opentelemetry} tails the access log and lands per-request metrics in
TimescaleDB \cite{timescaledb}; and a Redis \cite{carlson2013redis} control bus
carries decision envelopes between services. Four decision services (an anomaly
detector, a demand forecaster, a learned routing engine, and an autoscaler) run
off the request path, alongside a policy manager and an operator user interface.

The intelligence is built with the standard Python scientific stack: scikit-learn
and an Isolation Forest model \cite{liu2008isolationforest} for anomaly detection,
statsmodels with an ARIMA model \cite{box2015timeseries} for forecasting, and
Stable-Baselines3 \cite{raffin2021sb3} with a PPO policy \cite{schulman2017ppo}
for learned routing. Operational visibility comes through Prometheus
\cite{prometheus} and Grafana \cite{grafana}, and load is exercised with Locust
\cite{locust}.

The main result is that the closed control loop works end to end, and the
deterministic-fallback design keeps the whole system safe while the learned
components continue to improve.

Closed-loop proactive autoscaling is demonstrated. In a committed benchmark run,
the forecast-to-autoscaler-to-load-balancer loop scaled the backend pool from one
instance to six under a forecast-driven burst and then shrank it back as demand
subsided, across roughly twelve scaling actions, with per-phase p95 latencies
measured throughout.

For forecasting, the production default is a reliable moving-average predictor
that drives this autoscaling loop, with a trained ARIMA model available as an
optional advanced predictor under continued tuning. For anomaly detection, the
production default is a deterministic threshold detector chosen for guaranteed
stability; a trained Isolation Forest model agrees with this rule about 91.4\% of
the time while catching extra variance, and ships as an enhanced opt-in mode.

The learned PPO router performs on par with the strong, industry-standard
round-robin baseline in shadow mode. The deterministic fallback and operator
override keep routing at least as safe and effective as the established approach,
and richer gains on highly varied workloads are an active future direction.

Across all of these, the system runs safely on its deterministic baseline today,
and each learned component is a measurable, named enhancement on top of that
foundation. The full quantitative results are reported in \cref{ch:discussion}.

\section{Report Organization}
\label{sec:intro:organization}

The rest of this report is organised as follows. \Cref{ch:litreview} (Literature
Review) surveys the prior work this project builds on and positions against:
classical and learned load balancing, reactive versus proactive autoscaling,
time-series demand forecasting, anomaly detection for service health, and the
autonomic-computing tradition that frames adaptive systems as closed control
loops.

\Cref{ch:design_problem} and \cref{ch:design_solution} (Proposed Design) state
the engineering problem formally and then present the SmartLoad solution: the
three-plane separation of data, control, and decision; the contracts that let
independently developed services cooperate; the routing safety hierarchy; and the
deterministic-fallback discipline that runs through every learned component.

\Cref{ch:impl_overview}, \cref{ch:impl_decision}, and \cref{ch:impl_testing}
(Implementation) describe how the design was realised: the containerised stack
and telemetry pipeline, the internals of the four decision engines and their
training, and the testing and benchmarking strategy used to validate the system
end to end.

\Cref{ch:discussion} (Discussion and Project Impact) reports and interprets the
results: what was demonstrated, how each learned component compares with its
deterministic baseline, and the project's impact across the economic, societal,
environmental, and cultural contexts introduced in this chapter.

\Cref{ch:conclusion} (Conclusion) summarises the contribution, restates the
central idea that routing may learn while safety does not, and sets out the
Phase~2 directions (multi-tenancy, access control, and a managed control plane)
that follow from the single-tenant foundation built here.
